Since $n\ge2$, $m/2=2^{n-2}$ is an integer and $\Gamma(2+m/2)=(1+m/2)!=k!$ with $k=1+m/2\ge2$ (for a positive integer $k$, $\Gamma(k+1)=k!$). Robbins' form of Stirling's formula (\cite{Rob55}) gives $k!\le\sqrt{2\pi k}\,(k/e)^k e^{1/(12k)}$, i.e.\ $\log k!\le\frac12\log(2\pi k)+k(\log k-1)+\frac1{12k}$. Taking logarithms in the definition of $\widehat C_n$ (all factors are positive: $U(q)>0$, $\log(2+\sqrt5)>0$),
\begin{align*} \log\widehat C_n&=\log2+\frac m2\log\frac4\pi+\log k!+m\log\frac{U(q)}{\log(2+\sqrt5)}\\ &\le\log2+\frac m2\log\frac4\pi+\frac12\log(2\pi k)+k(\log k-1)+\frac1{12k}\\ &\qquad+m\log\frac{U(q)}{\log(2+\sqrt5)}\;=\;T_n . \end{align*}
If $d_n(\ell)\log\ell>T_n$ then $\ell^{d_n(\ell)}=e^{d_n(\ell)\log\ell}>e^{T_n}\ge\widehat C_n$, and Corollary~\ref{cor:Ahat} gives $\ell\nmid k_n$.

\emph{Check at $n=7$.} $m=64$, $k=33$, $q=512$, $U(512)=2.2431\ldots$: $T_7=\log2+32\log(4/\pi)+\frac12\log(66\pi)+33(\log33-1)+\frac1{396}+64\log(2.24312/1.44364)=0.69315+7.73006+2.66719+82.38475+0.00253+28.20496=121.68264$, while $\log\widehat C_7=121.68264$ ($\widehat C_7=7.016\times10^{52}$, \path{sage/r16_hatCn/r16_hatCn_pilot.log}); $T_7-\log\widehat C_7=7.7\times10^{-8}$, i.e.\ $1/(12k)-r_{33}$ with $r_{33}$ the Robbins remainder ($1/397<r_{33}<1/396$). Every constant is displayed in $T_n$; the remark that $T_n=\frac m2\log m+m\log(n+2)+O(m)$ (from $\log k!=k\log k-k+O(\log k)$ and $U(2^{n+2})=\frac{(n+2)\log2}{4}+O(1)$) is descriptive and is not part of the statement.
